A massive matter unit is a two-twistor composite whose flat mass is the bracket, $m^{2} = 2|\br{\pi^{1}\pi^{2}}|^{2}$ at $\Lambda = 0$. A fundamental matter unit is built of up to three twistors, each the vacuum of three pairs; the three-twistor composite is the maximal one, and the Standard-Model paper states the bound at observation level, which is the level it has here. The mass bracket transforms by $\det(U)$ under the internal $U(2)$, an $SU(2)$ singlet with charge one under the determinant character, its modulus the mass and its phase pure gauge. Under the global form $S(U(2) \times U(3))$ with the hypercharge of the Standard-Model paper it transforms as $(1,1)$ with $Y = +1$: the quantum numbers of $e^{c}$, and not those of a weak doublet. The relation between this object and the Standard Model's order parameter, an $SU(2)$ doublet of hypercharge one half, is not settled in the series, and the reading of the bracket as an order parameter is recorded at observation level only. \emph{Provenance:} \cite[Prop.~4.1]{ButlerT11}; \cite[Rem.~3.3]{ButlerT12} for the three-twistor bound, at the observation level that remark declares (Thm.~4.1 there presupposes it in naming the maximal hierarchy); \cite[Props.~3.1-3.2]{ButlerT12} for the transformation law and the quantum numbers, with the level of the order-parameter sentence as stated there.
